\appendix
\section{Planes mensuales: pedido y llegada no son intercambiables}\label{anexo:planes}

La primera tabla recoge \textbf{pedidos del plan inicial decididos durante A013}.
Enero corresponde a llegada en abril; diciembre, a marzo de A014. Son cantidades
enteras de una propuesta inicial, no medias ni medianas de decisiones revisadas.
La segunda tabla recoge las \textbf{llegadas iniciales previstas en A013}.
Enero y febrero pertenecen a compromisos anteriores estimados.

\begin{table}[htbp]\centering
\caption{Pedidos iniciales por mes de decisión A013, coches.}\label{tab:pedidos}
{\small\setlength{\tabcolsep}{3pt}\begin{tabular}{lrrrrrrrrrrrr}
\toprule
Segmento & 1 & 2 & 3 & 4 & 5 & 6 & 7 & 8 & 9 & 10 & 11 & 12 \\
\midrule
B & 153 & 132 & 168 & 121 & 66 & 45 & 67 & 78 & 100 & 64 & 103 & 136 \\
C & 0 & 0 & 0 & 0 & 0 & 0 & 0 & 172 & 220 & 163 & 201 & 232 \\
MD & 0 & 15 & 59 & 36 & 17 & 29 & 42 & 33 & 33 & 31 & 40 & 45 \\
SSUV & 659 & 494 & 750 & 530 & 253 & 245 & 325 & 412 & 496 & 219 & 335 & 426 \\
CSUV & 0 & 218 & 347 & 310 & 149 & 112 & 191 & 253 & 272 & 122 & 185 & 257 \\
MSUV & 353 & 622 & 693 & 593 & 293 & 336 & 452 & 516 & 633 & 314 & 309 & 468 \\
S & 80 & 72 & 74 & 30 & 0 & 0 & 0 & 0 & 0 & 56 & 83 & 6 \\
Total & 1.245 & 1.553 & 2.091 & 1.620 & 778 & 767 & 1.077 & 1.464 & 1.754 & 969 & 1.256 & 1.570 \\
\bottomrule
\end{tabular}
}
\fuente{\archivo{resultados/escenarios/tablas/plan_pedidos_inicial_a013_a014.csv};
filtro por año de decisión A013}
\end{table}

\begin{table}[htbp]\centering
\caption{Llegadas del plan inicial por mes de recepción A013, coches.}\label{tab:llegadas}
{\small\setlength{\tabcolsep}{3pt}\begin{tabular}{lrrrrrrrrrrrr}
\toprule
Segmento & 1 & 2 & 3 & 4 & 5 & 6 & 7 & 8 & 9 & 10 & 11 & 12 \\
\midrule
B & 124 & 124 & 81 & 153 & 132 & 168 & 121 & 66 & 45 & 67 & 78 & 100 \\
C & 454 & 454 & 0 & 0 & 0 & 0 & 0 & 0 & 0 & 0 & 172 & 220 \\
MD & 43 & 43 & 0 & 0 & 15 & 59 & 36 & 17 & 29 & 42 & 33 & 33 \\
SSUV & 390 & 390 & 1.135 & 659 & 494 & 750 & 530 & 253 & 245 & 325 & 412 & 496 \\
CSUV & 489 & 489 & 0 & 0 & 218 & 347 & 310 & 149 & 112 & 191 & 253 & 272 \\
MSUV & 576 & 576 & 0 & 353 & 622 & 693 & 593 & 293 & 336 & 452 & 516 & 633 \\
S & 34 & 34 & 329 & 80 & 72 & 74 & 30 & 0 & 0 & 0 & 0 & 0 \\
Total & 2.110 & 2.110 & 1.545 & 1.245 & 1.553 & 2.091 & 1.620 & 778 & 767 & 1.077 & 1.464 & 1.754 \\
\bottomrule
\end{tabular}
}
\fuente{\archivo{resultados/escenarios/tablas/plan_llegadas_inicial_a013.csv}}
\end{table}

Por ejemplo, los \num{1245} coches pedidos en enero de A013 aparecen como
llegadas de abril. Los \num{2110} previstos para enero ya vienen en camino
y no aparecen como un pedido que se pueda decidir en ese mismo enero.
La diferencia anual no es un fallo de balance: cada tabla utiliza un criterio
temporal diferente.
\FloatBarrier

\clearpage
\begin{landscape}
\pagestyle{plain}
\section{Tabla de resultados completa por segmento y escenario}\label{anexo:tabla5}
\noindent\small Medias de 1.000 réplicas. Los totales se calculan por réplica
antes del resumen; cobertura y servicio no son sumas de ratios. El redondeo puede
producir diferencias de una unidad. Los intereses de concesionarios se informan
aparte de la contribución de la marca.
\fuente{\archivo{resultados/escenarios/tablas/tabla5_plan_recomendado_a013.csv}}
{\footnotesize\renewcommand{\arraystretch}{1.02}\setlength{\tabcolsep}{5pt}\begin{longtable}{@{}llrrrrrrrr@{}}
\caption{Volúmenes A013: medias; cantidades en coches salvo cobertura.}\label{tab:anexo-volumenes}\\
\toprule
Seg. & Escenario & \makecell{Demanda\\(coches)} & \makecell{Ventas\\(coches)} & \makecell{Pedidos\\A013} & \makecell{Llegadas\\A013} & \makecell{Stock\\medio} & \makecell{Cobertura\\(meses)} & \makecell{Stock\\final} & \makecell{Stock $>90$\\al cierre} \\
\midrule
\endfirsthead
\toprule
Seg. & Escenario & \makecell{Demanda\\(coches)} & \makecell{Ventas\\(coches)} & \makecell{Pedidos\\A013} & \makecell{Llegadas\\A013} & \makecell{Stock\\medio} & \makecell{Cobertura\\(meses)} & \makecell{Stock\\final} & \makecell{Stock $>90$\\al cierre} \\
\midrule
\endhead
\midrule
\multicolumn{10}{r}{\footnotesize Continúa en la página siguiente.}\\
\endfoot
\bottomrule
\endlastfoot
B & Pesimista & 1.040 & 426 & 0 & 329 & 54,2 & 0,7 & 0 & 0 \\
C & Pesimista & 1.806 & 1.796 & 0 & 908 & 1.072,7 & 7,2 & 219 & 219 \\
MD & Pesimista & 360 & 360 & 286 & 283 & 90,9 & 3,1 & 48 & 5 \\
SSUV & Pesimista & 4.720 & 1.915 & 0 & 1.915 & 169,3 & 0,4 & 0 & 0 \\
CSUV & Pesimista & 2.838 & 2.816 & 1.751 & 2.325 & 979,5 & 4,3 & 520 & 245 \\
MSUV & Pesimista & 4.567 & 4.556 & 4.335 & 4.692 & 936,3 & 2,5 & 672 & 17 \\
S & Pesimista & 322 & 282 & 0 & 397 & 186,4 & 9,3 & 116 & 116 \\
\rowcolor{fondo}
Total & Pesimista & 15.652 & 12.151 & 6.371 & 10.848 & 3.489,4 & 2,7 & 1.574 & 601 \\
\addlinespace[5pt]
B & Central & 1.261 & 1.212 & 1.170 & 1.212 & 180,6 & 1,8 & 97 & 8 \\
C & Central & 2.267 & 2.258 & 959 & 1.298 & 985,2 & 5,3 & 147 & 11 \\
MD & Central & 435 & 433 & 375 & 349 & 83,7 & 2,3 & 41 & 1 \\
SSUV & Central & 5.668 & 5.387 & 4.882 & 5.867 & 700,8 & 1,5 & 480 & 22 \\
CSUV & Central & 3.394 & 3.327 & 2.357 & 2.791 & 919,8 & 3,4 & 475 & 138 \\
MSUV & Central & 5.601 & 5.554 & 5.501 & 5.588 & 862,6 & 1,9 & 570 & 5 \\
S & Central & 390 & 368 & 513 & 778 & 342,6 & 12,2 & 411 & 363 \\
\rowcolor{fondo}
Total & Central & 19.017 & 18.539 & 15.756 & 17.883 & 4.075,3 & 2,6 & 2.221 & 547 \\
\addlinespace[5pt]
B & Optimista & 1.930 & 1.578 & 1.602 & 1.540 & 136,8 & 0,9 & 58 & 1 \\
C & Optimista & 3.633 & 3.451 & 2.283 & 2.403 & 773,2 & 2,6 & 59 & 0 \\
MD & Optimista & 619 & 571 & 546 & 471 & 70,0 & 1,4 & 25 & 0 \\
SSUV & Optimista & 8.560 & 7.077 & 6.822 & 7.388 & 524,4 & 0,8 & 311 & 0 \\
CSUV & Optimista & 5.173 & 4.692 & 3.935 & 4.080 & 804,3 & 1,9 & 399 & 42 \\
MSUV & Optimista & 7.975 & 7.067 & 7.292 & 6.886 & 702,8 & 1,1 & 355 & 0 \\
S & Optimista & 556 & 516 & 788 & 964 & 346,3 & 8,7 & 449 & 361 \\
\rowcolor{fondo}
Total & Optimista & 28.445 & 24.952 & 23.268 & 23.731 & 3.357,8 & 1,4 & 1.656 & 404 \\
\addlinespace[5pt]
\end{longtable}
\clearpage
\begin{longtable}{@{}llrrrrrrrrrr@{}}
\caption{Economía A013: medias; importes en M\euro{} salvo cantidades de pérdidas y transfers.}\label{tab:anexo-economia}\\
\toprule
Seg. & Escenario & \makecell{Pérdidas\\(coches)} & \makecell{Transfers\\(coches)} & \makecell{Coste\\transfer} & \makecell{Coste\\stock} & VM & CVP & CO$_2$ & Arancel & Cambio & Resultado \\
\midrule
\endfirsthead
\toprule
Seg. & Escenario & \makecell{Pérdidas\\(coches)} & \makecell{Transfers\\(coches)} & \makecell{Coste\\transfer} & \makecell{Coste\\stock} & VM & CVP & CO$_2$ & Arancel & Cambio & Resultado \\
\midrule
\endhead
\midrule
\multicolumn{12}{r}{\footnotesize Continúa en la página siguiente.}\\
\endfoot
\bottomrule
\endlastfoot
B & Pesimista & 614 & 0 & 0,00 & 0,04 & 0,73 & 0,05 & 1,09 & 0,16 & 0,00 & -1,24 \\
C & Pesimista & 9 & 0 & 0,00 & 0,52 & 5,39 & 1,21 & 3,51 & 0,34 & 0,00 & -3,16 \\
MD & Pesimista & 0 & 0 & 0,00 & 0,06 & 1,34 & 1,38 & 0,75 & 0,12 & 0,00 & 0,45 \\
SSUV & Pesimista & 2.805 & 0 & 0,00 & 0,12 & 3,79 & 5,47 & 5,42 & 0,96 & 0,00 & -1,03 \\
CSUV & Pesimista & 22 & 0 & 0,00 & 0,68 & 6,90 & 10,48 & 5,38 & 0,90 & 0,00 & 3,52 \\
MSUV & Pesimista & 11 & 0 & 0,00 & 0,68 & 8,89 & 26,56 & 19,15 & 2,01 & 0,00 & 4,72 \\
S & Pesimista & 40 & 0 & 0,00 & 0,10 & 0,59 & 1,34 & 1,27 & 0,14 & 0,00 & -0,16 \\
\rowcolor{fondo}
Total & Pesimista & 3.501 & 0 & 0,00 & 2,20 & 27,63 & 46,49 & 36,56 & 4,64 & 0,00 & 3,09 \\
\addlinespace[5pt]
B & Central & 50 & 0 & 0,00 & 0,13 & 2,08 & 0,15 & 0,00 & 0,00 & 0,00 & 0,02 \\
C & Central & 9 & 0 & 0,00 & 0,55 & 6,60 & 1,69 & 0,00 & 0,00 & 0,00 & 1,15 \\
MD & Central & 2 & 0 & 0,00 & 0,06 & 1,59 & 1,69 & 0,00 & 0,00 & 0,00 & 1,63 \\
SSUV & Central & 281 & 0 & 0,00 & 0,51 & 10,66 & 15,37 & 0,00 & 0,00 & 0,00 & 14,86 \\
CSUV & Central & 67 & 0 & 0,00 & 0,66 & 8,11 & 12,43 & 0,00 & 0,00 & 0,00 & 11,77 \\
MSUV & Central & 47 & 0 & 0,00 & 0,63 & 10,68 & 32,54 & 0,00 & 0,00 & 0,00 & 31,91 \\
S & Central & 23 & 0 & 0,00 & 0,20 & 0,77 & 1,75 & 0,00 & 0,00 & 0,00 & 1,55 \\
\rowcolor{fondo}
Total & Central & 478 & 0 & 0,00 & 2,73 & 40,50 & 65,62 & 0,00 & 0,00 & 0,00 & 62,88 \\
\addlinespace[5pt]
B & Optimista & 352 & 0 & 0,00 & 0,10 & 2,71 & 0,20 & 0,00 & 0,00 & -0,37 & 0,47 \\
C & Optimista & 182 & 0 & 0,00 & 0,53 & 8,73 & 3,94 & 0,00 & 0,00 & -0,59 & 4,00 \\
MD & Optimista & 48 & 0 & 0,00 & 0,05 & 2,07 & 2,25 & 0,00 & 0,00 & -0,11 & 2,31 \\
SSUV & Optimista & 1.483 & 0 & 0,00 & 0,38 & 13,99 & 20,22 & 0,00 & 0,00 & -1,77 & 21,60 \\
CSUV & Optimista & 481 & 0 & 0,00 & 0,59 & 11,29 & 17,68 & 0,00 & 0,00 & -0,92 & 18,01 \\
MSUV & Optimista & 907 & 0 & 0,00 & 0,51 & 13,51 & 41,49 & 0,00 & 0,00 & -1,63 & 42,61 \\
S & Optimista & 40 & 0 & 0,00 & 0,23 & 1,07 & 2,47 & 0,00 & 0,00 & -0,13 & 2,37 \\
\rowcolor{fondo}
Total & Optimista & 3.493 & 0 & 0,00 & 2,38 & 53,36 & 88,23 & 0,00 & 0,00 & -5,52 & 91,37 \\
\addlinespace[5pt]
\end{longtable}
}
\end{landscape}
\pagestyle{fancy}
\clearpage

\section{Cómo reproducir y comprobar el trabajo}\label{anexo:reproducir}

\subsection{La cadena de ejecución}
Desde la raíz del repositorio se ejecutan, con un kernel que tenga las dependencias
instaladas, los cuatro notebooks en orden:
\begin{center}
\archivo{datos.ipynb} $\longrightarrow$ \archivo{demanda.ipynb}
$\longrightarrow$ \archivo{simulacion.ipynb}
$\longrightarrow$ \archivo{escenarios.ipynb}.
\end{center}
Cada uno conserva explicación, código y resultados. Las definiciones compartidas
son celdas \archivo{api}; no se requiere ejecutar manualmente celdas fuera de orden.
Las tablas y figuras se guardan en la carpeta de resultados de cada notebook.
La última fase regenera también el resumen y la presentación.

La ejecución de referencia utiliza semilla 42, 200 réplicas de piloto, 1.000
finales por escenario y 18 meses simulados. Las versiones registradas incluyen
Python 3.13.15, NumPy 2.5.3 y pandas 3.0.6; el registro completo está en
\archivo{resultados/escenarios/tablas/configuracion_ejecucion.json}.

\subsection{Qué verifican las comprobaciones}
\noindent\begin{tabularx}{\linewidth}{@{}p{3.1cm}Y@{}}\toprule
Bloque & Comprobaciones relevantes \\\midrule
Datos & Claves, fechas, lecturas de ausentes, identidad de clientes netos,
balances originales y ajuste de episodios aplicado una sola vez.\\
Previsión & Mercado desde A001; objetivo observado; CSUV fuera de validación;
mutar economía futura no altera la previsión operativa; mutar datos posteriores
no cambia el efecto de producto estimado al origen.\\
Simulador & Ejemplo de 286 coches; FIFO; coste de 15 días; límite de 180;
intereses por tramos; transfers con VM; unidades enteras y balances.\\
Integración & Mismo plan inicial; ausencia de anticipación; pedido de diciembre
a marzo; identidades económicas; totales por réplica y ledger absoluto.\\
Entrega & Coincidencia con tablas; documentos regenerables; legibilidad y
estructura de Word y PowerPoint.\\\bottomrule
\end{tabularx}

Comprobar una identidad contable no demuestra que la hipótesis de demanda sea
correcta. Por eso conviven controles de código, evaluación histórica y
sensibilidades metodológicas: cada uno responde a un tipo diferente de error.

\subsection{Qué representa este informe}
Las tablas y figuras de \archivo{docs/report/} son una instantánea de la ejecución
citada. Compilar LaTeX no vuelve a estimar los modelos. Las figuras vectoriales
didácticas se dibujan en el propio documento a partir de reglas o de copias de
resultados; no crean una segunda simulación. Si cambia el análisis, deben
actualizarse también las cifras y las interpretaciones del informe.

\section{Preguntas de repaso y errores de interpretación frecuentes}\label{anexo:repaso}

\begin{enumerate}[itemsep=7pt]
\item \textbf{¿Por qué no son iguales matriculaciones y demanda de clientes?}
Porque las primeras incluyen automatriculaciones. Además, ni siquiera las compras
netas identifican toda la demanda potencial cuando falta stock.
\item \textbf{¿Puede crecer el total y caer la gama comparable?}
Sí. CSUV incorpora clientes nuevos al perímetro de la marca aunque los segmentos
anteriores permanezcan estables o retrocedan.
\item \textbf{¿Por qué no se descuenta dos veces el adelanto?}
El histórico ajustado sirve para interpretar o normalizar el ancla; el registro
futuro se aplica una sola vez a A013. Replicar el descuento anual en cada mes
o añadir otro efecto de caída sobre la misma cantidad duplicaría su impacto.
\item \textbf{¿Un retraso económico de tres meses permite conocer todo A013?}
No. Sólo parte del horizonte queda cubierta por información ya observada; para
el resto se necesita congelación o escenarios.
\item \textbf{¿Qué significa MAE relativo del 14,8\,\%?}
Es el error absoluto mensual del total comparable en la evaluación retrospectiva
A012, dividido por sus ventas observadas. No valida CSUV ni garantiza el error A013.
\item \textbf{¿La cifra prevista es un pedido a fábrica?}
No. Antes se descuentan inventario y compromisos, se añade el colchón y se respeta
el plazo de llegada.
\item \textbf{¿95\,\% de servicio significa un 95\,\% de probabilidad de no tener roturas?}
No. El objetivo es la media de la proporción anual de unidades atendidas en el
central. Un segmento puede cumplirla y tener réplicas por debajo del umbral.
\item \textbf{¿Las medianas de segmentos suman la mediana del total?}
No necesariamente. Hay que sumar segmentos dentro de cada réplica y calcular
después el percentil del total.
\item \textbf{¿Toda venta perdida es una pérdida económica positiva?}
No. La venta adicional también habría soportado sus costes. Si la contribución
marginal es negativa, perder esa operación no equivale a perder un margen positivo.
\item \textbf{¿Más réplicas resuelven la incertidumbre de un arancel supuesto?}
No. Reducen ruido de Monte Carlo bajo ese arancel. Para estudiar otro tamaño
hay que cambiar el supuesto y comparar una sensibilidad.
\item \textbf{¿La mejor política evaluada es el óptimo global?}
No. Es la mejor dentro de la malla y del objetivo acordados, bajo las restricciones
y aproximaciones representadas.
\item \textbf{¿Qué dato vigilarías primero en A013?}
Ventas respecto a previsión, disponibilidad efectiva y compromisos; junto con
confirmación de costes, lanzamientos e interrupciones. Ninguna señal debe
interpretarse sin las restricciones que pueden estar ocultando demanda.
\end{enumerate}
